Este documento apresenta a versão 2 do modelo LaTeX comunitário para trabalhos acadêmicos da Universidade Federal do Ceará, desenvolvido sobre a classe abntexto-ufc e a infraestrutura abntexto. O trabalho tem como objetivo oferecer uma base editorial reproduzível para trabalhos de conclusão de curso, monografias de especialização, dissertações, teses e projetos de pesquisa, preservando requisitos institucionais compatíveis com as normas vigentes da Associação Brasileira de Normas Técnicas. A implementação organiza metadados, elementos pré-textuais, estrutura textual em seções, citações, referências, ilustrações, tabelas, quadros, códigos, algoritmos e elementos pós-textuais por meio de interfaces públicas e módulos opcionais. A validação combina testes semânticos e geométricos em PDF, matriz de perfis, compilação com pdfLaTeX e LuaLaTeX, Biber, glossários, índice e verificações de compatibilidade. Para o fluxo de depósito institucional, o documento de referência declara conformidade PDF/A-2b e a validação de release utiliza o veraPDF. Os resultados da suíte demonstram um fluxo consistente para produção, revisão e manutenção do modelo, sem substituir a conferência final das exigências específicas do curso, programa ou edital aplicável.

\palavraschave{LaTeX; normalização; trabalhos acadêmicos; UFC.}
